% Living Showcase: Verfahren und Listen (Palette Slot 2).
\StartChapter[ch:lists]{Verfahren und Listen}

\SubsectionBar[sec:procedure]{Nummerierte Schrittfolge}

\begin{ZSFlist}[Schritte mit Marke und Text\ZSFTitleTag{ordered}][ordered]
  \ZSFItem{Automatisch gezählt}{Die Nummerierung kommt aus der Umgebung.}
  \ZSFItem{Zweiteiliger Schritt}{Fette Marke oben, Erklärung darunter.}
  \ZSFItem{Nur bei stabiler Folge}{Sonst ist eine Liste die bessere Wahl.}
\end{ZSFlist}

\SubsectionBar[sec:lists]{Faktenliste ohne und mit Rahmen}

\begin{ZSFlist}
  \ZSFItem{ZSFlist}{Ohne Titelargument rahmenlos und ohne Kopfzeile.}
  \ZSFItem{ZSFItem}{Jeder Eintrag ist ein Paar aus Marke und Erklärung.}
  \ZSFItem{Farbiger Punkt}{Das Aufzählungszeichen trägt die Kapitelfarbe.}
  \ZSFItem{Die Marke ist optional — ein Eintrag ohne sie steht als blosse
    Aussage da, ohne Trenner, der ins Leere zeigt.}
\end{ZSFlist}

\begin{ZSFlist}[Dieselbe Liste mit Titel\ZSFTitleTag{mit Titel}]
  \ZSFItem{Ein Argument}{Der Titel bindet die Liste in eine leise Box.}
  \ZSFItem{Gleiche Einträge}{Marke und Aufzählungszeichen bleiben identisch.}
\end{ZSFlist}

\begin{runintext}
  Eine Umgebung, vier Erscheinungsformen: \ZSFkeyword{ZSFlist} entscheidet
  über Titel und \ZSFkeyword*{ordered}, nicht über den Namen.
\end{runintext}
